\paragraph{RQ 01 (frequência de deploy).}
Esperamos que a maioria dos repositórios populares com GitHub Actions publique releases com frequência semanal ou menor, e que apenas uma minoria alcance a faixa \textit{Elite} (deploys sob demanda, várias vezes por semana). Muitos projetos open-source seguem ciclos de release manuais ou por marcos, e a popularidade em estrelas indica visibilidade, não cadência de entrega; a distribuição deve ser assimétrica, com poucos repositórios muito frequentes e uma cauda longa de projetos que publicam poucas vezes ao ano.

\paragraph{RQ 02 (lead time).}
Esperamos lead times medianos de dias a poucas semanas, com a maioria dos repositórios nas faixas \textit{Medium} e \textit{Low}. Como o deploy é a publicação de uma release, e releases em projetos open-source acumulam commits por semanas antes de serem cortadas, o tempo entre o commit e a release tende a ser longo. Na variante (a), esperamos valores menores, porque ela é puxada pelos commits mais recentes antes de cada release; na variante (b), esperamos valores maiores e mais dispersos, porque ela incorpora os commits mais antigos de cada intervalo. Em ambas, a distribuição deve ter cauda longa e a mediana deve ser mais informativa que a média.

\paragraph{RQ 03 (taxa de falha de mudanças, CFR).}
Esperamos CFR baixo para a maior parte dos repositórios, na faixa \textit{Elite} ou \textit{High} (até cerca de 15\%), mas com grande dispersão. Na variante (a), baseada em workflow runs com falha, esperamos valores mais altos, porque falhas de CI ocorrem antes do deploy e incluem testes instáveis e problemas de infraestrutura que nunca chegam ao usuário. Na variante (b), baseada em releases corretivas, esperamos valores menores, porque só mudanças que chegaram a ser publicadas e depois corrigidas entram na conta, e a heurística tende a perder correções feitas sem nova release.

\paragraph{RQ 04 (tempo de recuperação).}
Esperamos que a maior parte das falhas seja recuperada em menos de um dia, já que o CI é refeito a cada push e corrigir uma falha costuma exigir um commit pequeno, mas com cauda longa de episódios que duram dias ou semanas em projetos de baixa atividade ou mantidos por poucas pessoas. Por isso, esperamos mediana na faixa \textit{Elite} ou \textit{High} (menos de um dia ou até uma semana) e média bem acima da mediana, além de uma parcela de episódios censurados, sem recuperação dentro da janela.
